% Options for packages loaded elsewhere
\PassOptionsToPackage{unicode}{hyperref}
\PassOptionsToPackage{hyphens}{url}
%
\documentclass[
]{article}
\usepackage{amsmath,amssymb}
\usepackage{iftex}
\ifPDFTeX
  \usepackage[T1]{fontenc}
  \usepackage[utf8]{inputenc}
  \usepackage{textcomp} % provide euro and other symbols
\else % if luatex or xetex
  \usepackage{unicode-math} % this also loads fontspec
  \defaultfontfeatures{Scale=MatchLowercase}
  \defaultfontfeatures[\rmfamily]{Ligatures=TeX,Scale=1}
\fi
\usepackage{lmodern}
\ifPDFTeX\else
  % xetex/luatex font selection
\fi
% Use upquote if available, for straight quotes in verbatim environments
\IfFileExists{upquote.sty}{\usepackage{upquote}}{}
\IfFileExists{microtype.sty}{% use microtype if available
  \usepackage[]{microtype}
  \UseMicrotypeSet[protrusion]{basicmath} % disable protrusion for tt fonts
}{}
\makeatletter
\@ifundefined{KOMAClassName}{% if non-KOMA class
  \IfFileExists{parskip.sty}{%
    \usepackage{parskip}
  }{% else
    \setlength{\parindent}{0pt}
    \setlength{\parskip}{6pt plus 2pt minus 1pt}}
}{% if KOMA class
  \KOMAoptions{parskip=half}}
\makeatother
\usepackage{xcolor}
\usepackage[margin=1in]{geometry}
\usepackage{graphicx}
\makeatletter
\def\maxwidth{\ifdim\Gin@nat@width>\linewidth\linewidth\else\Gin@nat@width\fi}
\def\maxheight{\ifdim\Gin@nat@height>\textheight\textheight\else\Gin@nat@height\fi}
\makeatother
% Scale images if necessary, so that they will not overflow the page
% margins by default, and it is still possible to overwrite the defaults
% using explicit options in \includegraphics[width, height, ...]{}
\setkeys{Gin}{width=\maxwidth,height=\maxheight,keepaspectratio}
% Set default figure placement to htbp
\makeatletter
\def\fps@figure{htbp}
\makeatother
\setlength{\emergencystretch}{3em} % prevent overfull lines
\providecommand{\tightlist}{%
  \setlength{\itemsep}{0pt}\setlength{\parskip}{0pt}}
\setcounter{secnumdepth}{-\maxdimen} % remove section numbering
\ifLuaTeX
  \usepackage{selnolig}  % disable illegal ligatures
\fi
\IfFileExists{bookmark.sty}{\usepackage{bookmark}}{\usepackage{hyperref}}
\IfFileExists{xurl.sty}{\usepackage{xurl}}{} % add URL line breaks if available
\urlstyle{same}
\hypersetup{
  pdfauthor={Parker Howell},
  hidelinks,
  pdfcreator={LaTeX via pandoc}}

\author{Parker Howell}
\date{Last Updated: 2023-09-25}

\begin{document}

\hypertarget{micro-i-notes}{%
\section{Micro I Notes}\label{micro-i-notes}}

\hypertarget{preference}{%
\subsection{Preference}\label{preference}}

\hypertarget{def-binary-relation}{%
\subsubsection{Def ``binary relation''}\label{def-binary-relation}}

A binary relation \(R\) on \(X\) is a subset of \(X \times X\).

The ordered pair \((x,y) \in R\) means \(x\) is (weakly) preferred to
\(y\).

\hypertarget{remark-about-notation}{%
\paragraph{Remark about notation}\label{remark-about-notation}}

We use \(\succsim\) to denote \(R\) but all these methods are the same
thing: \[(x,y)  \iff xRy \iff x \succsim y\]

\hypertarget{def-attributes-of-binary-relations}{%
\subsubsection{Def ``attributes of binary
relations''}\label{def-attributes-of-binary-relations}}

A binary relation \(\succsim\) is

\begin{enumerate}
\def\labelenumi{\arabic{enumi}.}
\tightlist
\item
  complete if for all \(x,y\in X\) either \(x \succsim y\) or
  \(y \succsim x\)
\item
  transitive if for all
  \(x,y,z\in X, [x \succsim y, y \succsim z] \Rightarrow x \succsim y\)
\item
  reflexive if for all \(x \in X, x \succsim x\)
\item
  symmetric if for all \(x,y\in X, x \succsim y \iff y \succsim x\)
\item
  irreflexive if for all \(x \in X, x \not\succsim x\)
\item
  asymmetric if for all
  \(x,y\in X, x \succsim y \Rightarrow y \not\succsim x\)
\end{enumerate}

\hypertarget{def-rational}{%
\subsubsection{Def ``rational''}\label{def-rational}}

A binary relation \(\succsim\) is \emph{rational} if it is complete and
transitive

\hypertarget{choice}{%
\subsection{Choice}\label{choice}}

\hypertarget{def-choice-function}{%
\subsubsection{Def ``choice function''}\label{def-choice-function}}

A function \(c: \mathcal{B} \rightarrow \mathcal{P}(X)\) is called a
\emph{choice function} if for all
\(A \in \mathcal{B}, c(A) \subseteq A\).

\begin{enumerate}
\def\labelenumi{\arabic{enumi}.}
\tightlist
\item
  If \(c\) is a choice function, then
  \(c(A) \neq \emptyset, \forall A \in \mathcal{B}\)
\item
  \(|c(A)| \geq 1\)
\end{enumerate}

\hypertarget{def-choice-structure}{%
\subsubsection{Def ``choice structure''}\label{def-choice-structure}}

The pair \((\mathcal{B}, c)\) is called a \emph{choice structure}.

\hypertarget{def-warp}{%
\subsubsection{Def ``WARP''}\label{def-warp}}

Weak Axiom of Revealed Preference. The following are equivalent for all
\(A, B \in \mathcal{B}\):

\begin{itemize}
\tightlist
\item
  \(c(A) \cap B \neq \emptyset \Rightarrow c(B) \cap \subseteq c(A)\)
\item
  \([c(A) \cap B \neq \emptyset \text{ AND } c(B) \cap A \neq \emptyset] \Rightarrow [c(A) \cap B \subseteq c(B) \text{ AND } c(B) \cap A \subseteq c(A)]\)
\item
  \(x, y \in A \cap B, x \in c(A), y \in c(B) \Rightarrow x \in c(B), y \in c(A)\)
\item
  Sen's \(\alpha\) and \(\beta\) both hold
\end{itemize}

\hypertarget{def-sens-alpha-and-beta}{%
\subsubsection{\texorpdfstring{Def Sen's \(\alpha\) and
\(\beta\)}{Def Sen's \textbackslash alpha and \textbackslash beta}}\label{def-sens-alpha-and-beta}}

For all \(A, B \in \mathcal{B}\):

\begin{enumerate}
\def\labelenumi{\arabic{enumi}.}
\tightlist
\item
  Sen's \(\alpha\): If \(x \in B \subseteq A\) and \(x \in c(A)\) then
  \(x\in c(B)\)
\item
  Sen's \(\beta\): If \(x,y \in c(B), B \subseteq A\) and \(y \in c(A)\)
  then \(x\in c(A)\)
\end{enumerate}

\hypertarget{prop-1-warp-iff-alpha-and-beta}{%
\subsubsection{\texorpdfstring{Prop 1 ``WARP \(\iff \alpha\) and
\(\beta\)''}{Prop 1 ``WARP \textbackslash iff \textbackslash alpha and \textbackslash beta''}}\label{prop-1-warp-iff-alpha-and-beta}}

\hypertarget{term-single-valued}{%
\paragraph{Term ``single-valued''}\label{term-single-valued}}

We say \(c\) is single-valued if
\(\forall A \in \mathcal{B}, |c(A)| = 1\)

\hypertarget{prop-2-alpha-and-single-valued-rightarrow-warp}{%
\subsubsection{\texorpdfstring{Prop 2 ``\(\alpha\) and single-valued
\(\Rightarrow\)
WARP''}{Prop 2 ``\textbackslash alpha and single-valued \textbackslash Rightarrow WARP''}}\label{prop-2-alpha-and-single-valued-rightarrow-warp}}

\hypertarget{choice-and-preference}{%
\subsection{Choice and Preference}\label{choice-and-preference}}

\hypertarget{def-c_succsim}{%
\subsubsection{\texorpdfstring{Def
``\(c_\succsim\)''}{Def ``c\_\textbackslash succsim''}}\label{def-c_succsim}}

For any binary relation \(\succsim\) define
\[c_\succsim(A) = \{x \in A: x \succsim y, \forall y \in A\}, \forall A \in \mathcal{P}(X)\]

\hypertarget{def-succsim_c}{%
\subsubsection{\texorpdfstring{Def
``\(\succsim_c\)''}{Def ``\textbackslash succsim\_c''}}\label{def-succsim_c}}

For a choice structure \((\mathcal{P}(X), c)\) define a binary relation
\(\succsim_c\) such that \[x \succsim_c y \iff x \in c(\{x,y\})\]

\hypertarget{prop-satisfies-warp-iff-succsim-rational}{%
\subsubsection{\texorpdfstring{Prop ``satisfies WARP \(\iff \succsim\)
rational}{Prop ``satisfies WARP \textbackslash iff \textbackslash succsim rational}}\label{prop-satisfies-warp-iff-succsim-rational}}

If \(\succsim\) is a binary relation on a finite set \(X\) then
\((\mathcal{P}(X), c_\succsim)\) is a choice structure that satisfies
WARP if and only if \(\succsim\) is rational.

\hypertarget{prop-warp-rightarrow-cc_succsim}{%
\subsubsection{\texorpdfstring{Prop ``WARP
\(\Rightarrow c=c_\succsim\)''}{Prop ``WARP \textbackslash Rightarrow c=c\_\textbackslash succsim''}}\label{prop-warp-rightarrow-cc_succsim}}

If \((\mathcal{P}, c)\) is a choice structure that satisfies WARP then
there exists a unique rational binary relation \(\succsim\) such that
\(c=c_\succsim\) (i.e., for all
\(A \in \mathcal{P}, c(A) = c_\succsim (A)\)).

\hypertarget{utility}{%
\subsection{Utility}\label{utility}}

\hypertarget{def-utility-function}{%
\subsubsection{Def ``utility function''}\label{def-utility-function}}

A function \(U: X \rightarrow \mathbb{R}\) represents \(\succsim\) if,
for all \(x, y \in X\) \[x \succsim y \iff U(x) \geq U(y).\]

\hypertarget{prop-u-represents-succsim-rightarrow-rational}{%
\subsubsection{\texorpdfstring{Prop ``\(U\) represents
\(\succsim \Rightarrow\)
rational''}{Prop ``U represents \textbackslash succsim \textbackslash Rightarrow rational''}}\label{prop-u-represents-succsim-rightarrow-rational}}

If \(U: X \rightarrow \mathbb{R}\) represents \(\succsim\), then
\(\succsim\) is rational.

Note: this proposition does NOT require finiteness!

\hypertarget{prop-x-finite-succsim-rational-exists-u}{%
\subsubsection{\texorpdfstring{Prop ``\(X\) finite + \(\succsim\)
rational
\Rightarrow \(\exists U\)}{Prop ``X finite + \textbackslash succsim rational \textbackslash exists U}}\label{prop-x-finite-succsim-rational-exists-u}}

If \(X\) is finite and \(\succsim\) is rational then there exists
\(U: X \rightarrow \mathbb{R}\) that represents \(\succsim\).

\hypertarget{prop-v-iff-phi-and-phi-strictly-increasing}{%
\subsubsection{\texorpdfstring{Prop ``\(V \iff \phi\) and \(\phi\)
strictly
increasing''}{Prop ``V \textbackslash iff \textbackslash phi and \textbackslash phi strictly increasing''}}\label{prop-v-iff-phi-and-phi-strictly-increasing}}

Suppose \(U: X \rightarrow \mathbb{R}\) represents \(\succsim\). Then
\(V: X \rightarrow \mathbb{R}\) represents \(\succsim\) if and only if
there exists \(\phi: U(X) \rightarrow V(X)\) such that
\(V = \phi \circ U\) and \(\phi\) is strictly increasing.

\hypertarget{expected-utility}{%
\subsection{Expected Utility}\label{expected-utility}}

\hypertarget{def-lottery}{%
\subsubsection{Def ``lottery''}\label{def-lottery}}

Suppose \(X\) is a non-empty and finite set of prizes (consequences).
Then a \emph{lottery}, \(p\) is a probability distribution over \(X\):
\[p: X \rightarrow [0,1], \text{ such that } \sum_{x\in X} p(x) = 1.\]

We denote \(\mathcal{L}\) as the set of all lotteries.

\hypertarget{term-degenerate-lottery}{%
\paragraph{Term ``degenerate lottery''}\label{term-degenerate-lottery}}

If \(p(x)=1\) for some \(x \in X\) then we write \(\delta_x\).

\hypertarget{remark-mathcall-is-the-set-of-alternatives-and-unlike-x-is-infinite}{%
\paragraph{\texorpdfstring{Remark: \(\mathcal{L}\) is the set of
alternatives, and unlike \(X\) is
infinite!}{Remark: \textbackslash mathcal\{L\} is the set of alternatives, and unlike X is infinite!}}\label{remark-mathcall-is-the-set-of-alternatives-and-unlike-x-is-infinite}}

\hypertarget{def-expected-utility-functionutility-index}{%
\subsubsection{Def ``expected utility function/utility
index''}\label{def-expected-utility-functionutility-index}}

\(U: \mathcal{L} \rightarrow \mathbb{R}\) is an expected utility
function if there exists \(u: X \rightarrow \mathbb{R}\) such that
\(U(p) = \sum_{x \in X} u(x)p(x), \forall p \in \mathcal{L}\).

Called a ``von Neumann--Morgenstern utility index'' or a ``Bernoulli
index.''

\hypertarget{def-mixture-operation}{%
\subsubsection{Def ``Mixture operation''}\label{def-mixture-operation}}

For all \(p, q \in \mathcal{L}, \alpha \in [0,1]\):

\begin{itemize}
\tightlist
\item
  let \(\alpha p + (1-\alpha)q\) denote a lottery \(r\in \mathcal{L}\)
  such that for all \(x\in X, r(x) = \alpha p(x) + (1-\alpha)q(x)\).
\end{itemize}

\hypertarget{def-expected-utility-representation}{%
\subsubsection{Def ``expected utility
representation''}\label{def-expected-utility-representation}}

\(\succsim\) on \(\mathcal{L}\) has an expected utility representation
if there exists \(u: X \rightarrow \mathbb{R}\) such that for
\(p,q \in \mathcal{L}\)
\[p \succsim q \iff U(p) = \sum_{x\in X} p(x)u(x) \geq U(q) = \sum_{x\in X} q(x)u(x)\]

\hypertarget{def-linear-function}{%
\subsubsection{Def ``linear function''}\label{def-linear-function}}

A function \(U: \mathcal{L} \rightarrow \mathbb{R}\) is \emph{linear} if
for all \(\alpha \in [0,1]\) and \(p,q \in \mathcal{L}\):
\[U(\alpha p + (1-\alpha)q = \alpha U(p) + (1-\alpha) U(q).\]

\hypertarget{lemma-linear-iff-u-is-an-expected-utility-function}{%
\subsubsection{\texorpdfstring{Lemma ``linear \(\iff U\) is an expected
utility
function''}{Lemma ``linear \textbackslash iff U is an expected utility function''}}\label{lemma-linear-iff-u-is-an-expected-utility-function}}

A function \(U: \mathcal{L} \rightarrow \mathbb{R}\) is \emph{linear} if
and only if \(U\) is an expected utility function.

\hypertarget{def-vnm-axioms}{%
\subsubsection{Def ``vNM Axioms''}\label{def-vnm-axioms}}

The von Neumann--Morgenstern axioms are:

\begin{itemize}
\tightlist
\item
  A1. (Rationality): \(\succsim\) is complete and transitive.
\item
  A2. (Independence): For all
  \(p,q,r \in \mathcal{L}, \alpha \in (0,1)\) we have
  \([p \succsim q] \Rightarrow [\alpha p + (1-\alpha)r \succsim \alpha q + (1-\alpha)r]\).
\item
  A3. (Continuity): For all \(p, q, r \in \mathcal{L}\), if
  \(p \succ q \succ r\) then there exists \(\alpha, \beta \in (0,1)\)
  such that \(\alpha p + (1- \alpha) r \succ q\) and
  \(q \succ \beta p + (1-\beta) r\).
\item
  We don't use A2* in this class, but some literature uses it:

  \begin{itemize}
  \tightlist
  \item
    A2*. (Independence, strengthened):
    \([p \succsim q] \iff [\alpha p + (1-\alpha) r \succsim \alpha q + (1-\alpha)r]\)
  \end{itemize}
\end{itemize}

\hypertarget{lemma-applying-vnm-axioms}{%
\paragraph{Lemma ``applying vNM
Axioms''}\label{lemma-applying-vnm-axioms}}

Suppose \(\succsim\) satisfies A1, A2, and A3 of vNM axioms.

\begin{enumerate}
\def\labelenumi{(\roman{enumi})}
\tightlist
\item
  \(\forall \alpha \in (0,1), p,q \in \mathcal{L}\):

  \begin{enumerate}
  \def\labelenumii{(\alph{enumii})}
  \tightlist
  \item
    \(p \succ q \Rightarrow p \succ \alpha p + (1 - \alpha) q \succ q\)
  \item
    \(p \sim q \Rightarrow p \sim \alpha p + (1 - \alpha) q \sim q\)
  \end{enumerate}
\item
  \(p \succ q, 1 \geq \alpha > \beta \geq 0 \Rightarrow \alpha p + (1-\alpha)q \succ \beta p + (1-\beta)q\)
\item
  \(p \sim q, \alpha \in (0,1) \Rightarrow \alpha p + (1-\alpha)r \sim \alpha q + (1-\alpha)r\)
\item
  \(p \succsim q, r \succsim s, \alpha \in (0,1) \Rightarrow \alpha p + (1-\alpha)r \succsim \alpha q + (1-\alpha)s\)
\end{enumerate}

\hypertarget{thm-a1-a3-iff-succsim-has-eu}{%
\subsubsection{\texorpdfstring{Thm ``A1-A3 \(\iff \succsim\) has
EU''}{Thm ``A1-A3 \textbackslash iff \textbackslash succsim has EU''}}\label{thm-a1-a3-iff-succsim-has-eu}}

Suppose \(X\) is finite and \(\succsim\) is a binary relation on
\(\mathcal{L}\). \(\succsim\) satisfies A1, A2, and A3 if and only if
\(\succsim\) has an Expected Utility representation.

\end{document}
